%%%%%%%%%%%%%%%%%%%%%%%%%%%%%%%%%%%%%%%%%%%%%%%%%%%%%%%%%%%%%%%%%%%%%%%%%%%%%%%%%%%

%\blankpage
\pagebreak
\thispagestyle{empty}
\addcontentsline{toc}{part}{\bfseries Marrocos}
\mbox{}
\begin{center}

\vspace{3cm}

{\brabo\HUGE Marrocos}
\end{center}
\pagebreak

%%%%%%%%%%%%%%%%%%%%%%%%%%%%%%%%%%%%%%%%%%%%%%%%%%%%%%%%%%%%%%%%%%%%%%%%%%%%%%%%%%%

\blankpage
\pagebreak
\thispagestyle{empty}
\addcontentsline{toc}{part}{Marraquexe}
\mbox{}
\begin{center}

\vspace{3cm}

% {\formularlight\small PARTE 1}

% \medskip

{\brabo\HUGE\itshape Marraquexe}
\end{center}
\pagebreak

%\part[Drama barroco e tragédia]{Drama barroco\\ e tragédia}

%%%%%%%%%%%%%%%%%%%%%%%%%%%%%%%%%%%%%%%%%%%%%%%%%%%%%%%%%%%%%%%%%%%%%%%%%%%%%%%%%%%

\chapter*{Jardim Majorelle}

Famoso jardim botânico e espaço cultural. Criado pelo artista francês Jacques Majorelle e posteriormente restaurado por Yves Saint Laurent e Pierre Bergé, tornou-se um dos pontos turísticos mais visitados da cidade, conhecido por sua vegetação exuberante, arquitetura marcante e o icônico azul \mbox{Majorelle}. 

Além dos jardins, o complexo abriga o Museu Berbere e áreas dedicadas ao legado de Yves Saint Laurent, tornando a visita interessante tanto para amantes de natureza quanto de arte e design. O site informa que o jardim abre diariamente, das 8h às 18h30, com última entrada às 18h. O museu Yves Saint Laurent tem programação própria e deve ser conferido separadamente.

\chapter{Le Jardin Secret}

Um dos jardins históricos mais interessantes da medina. O complexo foi restaurado a partir de um antigo \textit{riad} e permite entender como arquitetura, água e jardins estavam ligados à vida doméstica e palaciana de Marraquexe.

Há dois jardins principais, um de inspiração islâmica e outro mais exótico, além de uma torre de onde se vê a medina. É um bom contraponto ao movimento dos souks: vale entrar, diminuir o ritmo e observar a arquitetura hidráulica e a organização dos espaços. Em outubro, o horário oficial informado é das 9h30 às 18h30. A entrada na torre é cobrada separadamente.

\chapter*{Dar El Bacha (Musée\\ des Confluences)}
\addcontentsline{toc}{chapter}{Dar El Bacha (Musée des Confluences)}

Antigo palácio de Thami El Glaoui, Paxá de Marraquexe, construído no início do século \textsc{xx}. O palácio era residência, espaço de recepção e também cenário de exercício de poder durante o período do Protetorado Francês. Hoje abriga o Musée des Confluences Dar El Bacha. 

Mais do que uma coleção específica, o lugar vale pela própria arquitetura: pátios, \textit{zellij}, madeira de cedro, estuque, jardins e elementos que mostram o encontro entre tradições marroquinas e influências europeias. É um dos lugares em que arquitetura e história política se encontram de maneira mais clara.

\chapter{Koutoubia}

A mesquita Koutoubia é a maior mesquita e um dos monumentos mais representativos da cidade de Marraquexe. Seu minarete é o modelo dos das mesquitas de Rabate, com a Torre Haçane e de Sevilha com a Giralda. Construída pelos almóadas no século \textsc{xii}, sua torre tem cerca de 77 metros e domina a paisagem da cidade histórica. O minarete é particularmente importante para quem está interessado na relação entre Marrocos e Al-Andalus: sua linguagem arquitetônica tornou-se referência para a Torre Hassan, em Rabat, e para a Giralda, em Sevilha. A comparação é interessante porque mostra como formas arquitetônicas circularam pelo mundo islâmico ocidental.

A mesquita é um local de culto e o interior não é aberto a visitantes não muçulmanos; a visita, portanto, é sobretudo externa.

\chapter{Túmulos Saadianos} 

Mausoléu real da dinastia saadiana, que governou o Marrocos entre os séculos \textsc{xvi} e \textsc{xvii}. O conjunto está associado ao sultão Ahmad al-Mansur, cujo reinado marcou um dos períodos de maior esplendor da dinastia.

\chapter{Palácio el Badi}

O El Badi foi construído pelo sultão saadiano Ahmad al-Mansur a partir de 1578, depois da vitória marroquina na Batalha de Alcácer-Quibir. A riqueza da construção esteve ligada, entre outras coisas, ao resgate pago pelos portugueses após a \mbox{batalha.}

Hoje restam grandes muralhas, pátios e estruturas em ruínas. A graça da visita está justamente em imaginar a escala do palácio original e perceber a passagem do esplendor cortesão para a ruína.

\chapter{Palácio Bahia}

Construído no final do século \textsc{xix} para o grão-vizir Ba Ahmed, o Bahia é um grande complexo de pátios, jardins e aposentos. Seu nome costuma ser traduzido como ``brilho'' ou ``esplêndido''. Os jardins ocupam uma área de 8\,000 m² e as 150 divisões abrem-se para diversos pátios interiores.

É especialmente interessante para observar a combinação de \textit{zellij}, madeira de cedro, estuque e jardins. O próprio turismo oficial de Marraquexe destaca o trabalho de artesãos marroquinos e andaluzes na construção; materiais também chegaram de diferentes regiões do país.

\chapter{Mellah de Marraquexe}

É o antigo bairro judaico de Marraquexe. A sinagoga Salat Al Azama, também conhecida como Lazama, é um dos pontos centrais do bairro. Seu pequeno pátio, revestido de azulejos sobre a vida judaica, ajudam a compreender a história da comunidade. Há também o cemitério judaico de Miaara, um espaço particularmente impressionante. Vale observar as diferenças de escala, arquitetura e uso entre o Mellah e o restante da medina.

Nas proximidades, o Mercado Mellah concentra o comércio cotidiano de carnes e produtos agrícolas, enquanto o Bab Mellah Spice Souk reúne especiarias e outros produtos. Entre as ruas do bairro também se encontram cafés, restaurantes e pequenas barracas que servem carnes grelhadas e tajines.

\chapter{Jemaa el-Fna}

Jemaa el-Fna é a principal e mais célebre praça de Marraquexe. A praça faz parte da medina de Marraquexe, inscrita pela \textsc{unesco} na lista do Patrimônio Mundial. É um espaço de encontro e de expressão da cultura oral. Ao longo do dia, vendedores, músicos, artistas e comerciantes ocupam o espaço; ao anoitecer, a praça se transforma, com barracas de comida, músicos e pequenos círculos de contadores de histórias e artistas --- é especialmente interessante observar os grupos que se formam espontaneamente ao redor dos \textit{halqa}.

\chapter{Maison de la Photographie}

A Maison de la Photographie de Marrakech é um museu e espaço cultural dedicado à história da fotografia em Marrocos. Instalado em um riad tradicional da medina, reúne fotografias, cartões-postais, documentos e outros registros visuais que ajudam a reconstruir diferentes aspectos da vida marroquina entre o fim do século \textsc{xix} e meados do século \textsc{xx}.

A coleção inclui imagens de cidades, paisagens, arquitetura, festas, costumes e diferentes grupos da sociedade marroquina. Para além do interesse fotográfico, a visita oferece uma perspectiva rara sobre um Marrocos anterior às grandes transformações do século \textsc{xx} e permite observar lugares e modos de vida que hoje já não existem da mesma maneira.

O terraço do edifício também oferece uma das vistas mais bonitas sobre os telhados da medina e as montanhas do Atlas em dias de boa visibilidade. O ingresso da Maison de la Photographie também dá acesso ao Musée de la Musique de Marrakech.

\chapter{Souk Semmarine}

É o grande eixo comercial da medina e o mais turístico e movimentado. Isso não significa que seja ``falso'' --— ele é um souk histórico e continua vendendo artesanato marroquino --— mas é onde você encontrará mais facilmente tapetes, bolsas, roupas, babouches, luminárias, cerâmica, objetos para turistas. Para compras mais especializadas, consultar:

\begin{itemize}
\item Souk des Teinturiers: tecidos e tinturaria;
\item Souk Haddadine: metal, lanternas e objetos de ferro;
\item Souk El Attarine: objetos de cobre e metal;
\item Souk Cherratine: artigos de couro;
\item Rahba Kedima/\,Souk des Épices: especiarias e pequenos objetos;
\item Souk Smata: babouches e calçados;
\item Souk Zrabi: tapetes.
\end{itemize}

\chapter{Mercados de moedas}

O mercado é menos organizado do que o de artesanato. Antiguários e mercados de antiguidades podem ter peças interessantes, mas é importante distinguir objeto antigo, reprodução e peça de procedência incerta.

\begin{itemize}
\item Azdin Antiques: antiquário na medina; vale perguntar especificamente por monnaies anciennes marocaines.
\item Amazonite: antiquário em Gueliz.
\item The Moroccan Doors Souk el Khemis: antiquário/mercado de antiguidades, embora não seja especializado em numismática.
\end{itemize}

\chapter{Librairie Booklore} 

Livraria contemporânea em Guéliz, com seleção de ficção, arte, arquitetura, culinária e livros em francês e inglês. Para quem trabalha com livros, vale menos pela ideia de ``livraria turística'' e mais para observar o que uma livraria marroquina contemporânea escolhe colocar em destaque.

\chapter*{«Bouquinistes» de\\ Bab Doukkala}
\addcontentsline{toc}{chapter}{«Bouquinistes» de Bab Doukkala}

Para garimpo, esta é provavelmente a experiência mais divertida. São vendedores de livros usados na região de Bab Doukkala, com edições antigas e material que aparece de forma menos previsível do que em uma livraria organizada.

\chapter{Madrassa Ben Youssef}

A Madrassa Ben Youssef é uma histórica escola islâmica do século \textsc{xvi}, um dos grandes monumentos da arquitetura saadiana de Marraquexe. Funcionou como escola religiosa e centro de estudo, com quartos para estudantes organizados em torno de pátios e uma decoração exuberante de \textit{zellij}, estuque esculpido e madeira de cedro.

\chapter{Museu Al Maaden}

Museu dedicado à arte contemporânea africana. A visita muda completamente de registro em relação à medina e é uma boa maneira de lembrar que o Marrocos contemporâneo também participa de circuitos artísticos africanos mais amplos. Fica fora do centro histórico, portanto deve ser tratado como um passeio próprio, e não como uma parada rápida entre dois monumentos da medina. O site oficial informa funcionamento de quarta a domingo, das 10h às 18h.

\chapter{Comptoir des Mines Galerie}

Galeria contemporânea em Guéliz, interessante para entrar em contato com a produção artística atual do Marrocos e com a cena de arte contemporânea de Marraquexe.

\chapter{Musée Macma}

Museu privado de arte e cultura dedicado a artistas marroquinos históricos e contemporâneos. Pode funcionar como complemento ao Comptoir des Mines, sobretudo se vocês quiserem comparar diferentes maneiras de apresentar a arte marroquina \mbox{recente.}

\chapter*{Música andalusina no\\ Dar Zephyr}
\addcontentsline{toc}{chapter}{Música andalusina no Dar Zephyr}

Há uma programação anunciada para outubro de 2026 no Dar Zephyr, em um pátio da medina, com uma orquestra de doze músicos dedicada ao repertório andalusino. É uma opção especialmente alinhada ao interesse pela continuidade entre a tradição musical de Al-Andalus e o Marrocos.

A programação encontrada para 2026 aparece entre setembro e outubro; para os dias da viagem, é essencial conferir a data exata antes de comprar. A apresentação começa às 20h30, com portas às 20h. Os ingressos anunciados são de 400 \textsc{mad} para lugares no pátio e 600 \textsc{mad} para primeira fila; chá e doces estão incluídos na opção sentada. A programação deve ser reconfirmada para 10--12 de outubro de 2026 antes da reserva.

\chapter{Musée de la Musique}

Talvez seja a melhor opção do guia para ouvir música tradicional em um contexto relativamente íntimo. O museu funciona em uma antiga \textit{douiria} saadiana e apresenta a diversidade musical do Marrocos, incluindo música arabo-andalusina, amazigh e Gnawa.

O calendário publicado para setembro de 2026 a maio de 2027 informa concertos às segundas, quartas e sextas, às 18h: segunda-feira, música andalusina; quarta-feira, música amazigh; sexta-feira, Gnawa.
